% ==========================================================================
%  Chapter 70 — Higher-Order Prismatic Domains and VTK Local Axes
% ==========================================================================

\chapter{Higher-Order Prismatic Domains and VTK Local Axes}
\label{ch:ho-prismatic-domain}

Chapter~\ref{ch:reinforced-submodel} introduced the
\texttt{PrismaticDomainBuilder} for generating structured hexahedral
meshes (Hex8) automatically aligned to beam elements.  The present
chapter generalises the builder to support higher-order hexahedral
elements---specifically the 20-node serendipity element (Hex20) and the
27-node triquadratic element (Hex27)---and adds a new VTK cell-data
mechanism for exporting element local axes to ParaView.


% =========================================================================
\section{Motivation}
\label{sec:ho-prism-motivation}

Linear hexahedral elements (Hex8) are the simplest three-dimensional
continuum elements but suffer from well-known shear-locking artifacts
in bending-dominated problems.  When a prismatic sub-model is used to
evaluate nonlinear concrete behaviour after kinematic reconstruction
from a beam analysis, the accuracy of the strain field within the
sub-model can be significantly improved by using quadratic elements.

Two quadratic hexahedral variants are available in \texttt{fall\_n}:

\begin{description}
  \item[Hex20 — Serendipity (\texttt{SerendipityElement<3,3,2>})]
        Uses 8 corner nodes plus 12 edge-midpoint nodes (no face or
        body-centre nodes).  It yields a 20-DOF element that generally
        performs well in bending and avoids most Hex8 locking issues.
  \item[Hex27 — Triquadratic (\texttt{LagrangeElement3D<3,3,3>})]
        Uses all 27 nodes of a $3\times 3\times 3$ tensor-product grid.
        It offers full quadratic completeness at the cost of 7 additional
        internal nodes per element.
\end{description}

Both use $3\times 3\times 3$ Gauss--Legendre quadrature (27 integration
points), which is the minimum complete rule for triquadratic shape
functions.


% =========================================================================
\section{The \texttt{HexOrder} Enumeration}
\label{sec:ho-hex-order}

The generalisation is controlled through the \texttt{HexOrder} enumeration:

\begin{lstlisting}[language=C++, basicstyle=\ttfamily\small]
enum class HexOrder {
    Linear,       // Hex8  -- 8-node trilinear
    Serendipity,  // Hex20 -- 20-node serendipity
    Quadratic     // Hex27 -- 27-node triquadratic
};
\end{lstlisting}

A new field \texttt{hex\_order} was added to \texttt{PrismaticSpec} with
a default value of \texttt{HexOrder::Linear}, so every existing call
site remains unchanged:

\begin{lstlisting}[language=C++, basicstyle=\ttfamily\small]
struct PrismaticSpec {
    double width, height, length;
    int    nx, ny, nz;
    HexOrder hex_order = HexOrder::Linear;   // <-- new
    // ... origin, e_x, e_y, e_z, physical_group
};
\end{lstlisting}


% =========================================================================
\section{Structured Grid Generalisation}
\label{sec:ho-grid}

The \texttt{PrismaticGrid} metadata struct was rewritten to support
both linear and quadratic grids.  Two key additions drive the change:

\begin{enumerate}
  \item A \textbf{step} factor: \texttt{step = 1} for linear meshes,
        \texttt{step = 2} for quadratic meshes.
  \item The number of nodes per direction becomes
        \texttt{step $\cdot$ n + 1} instead of \texttt{n + 1}.
\end{enumerate}

\noindent For a quadratic mesh with $n_x \times n_y \times n_z$ elements
the grid has $(2n_x + 1) \times (2n_y + 1) \times (2n_z + 1)$ node
positions.  Node IDs are assigned in lexicographic (column-major) order:
%
\begin{equation}
  \text{node\_id}(i_x, i_y, i_z)
  = i_z \cdot N_x \cdot N_y + i_y \cdot N_x + i_x,
  \qquad
  N_\alpha = \text{step} \cdot n_\alpha + 1.
  \label{eq:ho-node-id}
\end{equation}

For Hex20 (serendipity) meshes, all $(2n+1)^3$ grid positions are
created as nodes in the \texttt{Domain}, even though face-centre and
body-centre positions are not referenced by any element connectivity.
These orphan nodes are harmless and greatly simplify the code.

The node spacing in each direction is $\Delta x / \text{step}$ so that
mid-edge nodes lie at exactly the half-way position between element
corners:

\begin{lstlisting}[language=C++, basicstyle=\ttfamily\small]
const double sub_dx = dx / static_cast<double>(step);
// lx = x0 + ix * sub_dx
\end{lstlisting}


% =========================================================================
\section{Element Connectivity Patterns}
\label{sec:ho-connectivity}

\subsection{Hex8 — Linear}

The original column-major pattern on a $\{0,1\}^3$ sub-grid:
%
\begin{equation}
  \texttt{conn}[i_0 + 2\,i_1 + 4\,i_2]
  = \text{node\_id}(e_x + i_0,\; e_y + i_1,\; e_z + i_2),
  \qquad i_\alpha \in \{0,1\}.
\end{equation}

\subsection{Hex20 — Serendipity}

The 20-node connectivity follows the node ordering of
\texttt{SerendipityCell<3,2>}:

\begin{itemize}
  \item Nodes 0--7: corners at even grid offsets
        $\{0,2\}^3$ from the base $(2e_x, 2e_y, 2e_z)$.
  \item Nodes 8--11: bottom-face edge midpoints (one odd $x$ or $y$
        index, $z$ at base).
  \item Nodes 12--15: top-face edge midpoints (one odd index,
        $z$ at base$+2$).
  \item Nodes 16--19: vertical edge midpoints (both $x,y$ even,
        $z$ at base$+1$).
\end{itemize}

\noindent This ordering matches VTK's
\texttt{VTK\_QUADRATIC\_HEXAHEDRON} exactly, so no permutation is
needed (identity map in \texttt{VTKCellTraits}).

\subsection{Hex27 — Triquadratic}

A full $3^3$ tensor product in column-major order:
%
\begin{equation}
  \texttt{conn}[i_0 + 3\,i_1 + 9\,i_2]
  = \text{node\_id}(2e_x + i_0,\; 2e_y + i_1,\; 2e_z + i_2),
  \qquad i_\alpha \in \{0,1,2\}.
\end{equation}


% =========================================================================
\section{Quadrature Selection}
\label{sec:ho-quadrature}

The quadrature rule is selected per element type:

\begin{center}
\begin{tabular}{lccc}
  \toprule
  Element & Template & Integrator & Gauss points \\
  \midrule
  Hex8  & \texttt{LagrangeElement3D<2,2,2>}
        & \texttt{GaussLegendreCellIntegrator<2,2,2>} & $2^3 = 8$ \\
  Hex20 & \texttt{SerendipityElement<3,3,2>}
        & \texttt{GaussLegendreCellIntegrator<3,3,3>} & $3^3 = 27$ \\
  Hex27 & \texttt{LagrangeElement3D<3,3,3>}
        & \texttt{GaussLegendreCellIntegrator<3,3,3>} & $3^3 = 27$ \\
  \bottomrule
\end{tabular}
\end{center}

The $3 \times 3 \times 3$ rule is the minimum complete rule for a
triquadratic integrand $(\text{order } 2 \times 2 = 4$,
requiring $\lceil 5/2 \rceil = 3$ points per direction).


% =========================================================================
\section{Reinforced Domain with Higher-Order Elements}
\label{sec:ho-reinforced}

The \texttt{make\_reinforced\_prismatic\_domain()} function was updated
to dispatch on \texttt{hex\_order} using the same logic as the
unreinforced builder.  The only change affecting rebar line elements is
that corner nodes in a quadratic grid are located at even indices:

\begin{lstlisting}[language=C++, basicstyle=\ttfamily\small]
PetscInt conn[2] = {
    grid.node_id(step * bar.ix, step * bar.iy, step * iz),
    grid.node_id(step * bar.ix, step * bar.iy, step * (iz + 1)),
};
\end{lstlisting}

Since \texttt{step = 2} for quadratic meshes, the rebar bar at
cross-section position \texttt{(ix, iy)} connects grid nodes
\texttt{(2·ix, 2·iy, 2·iz)} and \texttt{(2·ix, 2·iy, 2·(iz+1))},
both of which are corner positions shared by the surrounding hex
elements.  The perfect-bond assumption is thus preserved exactly.


% =========================================================================
\section{The \texttt{align\_to\_beam} Helper}
\label{sec:ho-align}

The convenience function \texttt{align\_to\_beam()} now accepts an
optional \texttt{HexOrder} parameter at the end of its argument list:

\begin{lstlisting}[language=C++, basicstyle=\ttfamily\small]
PrismaticSpec align_to_beam(
    A, B, up,
    width, height,
    nx, ny, nz,
    group = "Solid",
    hex_order = HexOrder::Linear    // <-- new, default = backward-compatible
);
\end{lstlisting}

The returned \texttt{PrismaticSpec} carries the requested
\texttt{hex\_order} alongside the computed origin and rotation matrix.


% =========================================================================
\section{VTK Local-Axes Cell Data}
\label{sec:ho-vtk-axes}

\subsection{Motivation}

When interpreting material state fields (stress, strain, crack
directions) on a prismatic sub-model, the analyst needs to know the
element's local coordinate frame relative to the global axes.  Without
this information, a shear stress $\sigma_{xz}$ in the results could
refer to any direction, making visual inspection unreliable.

\subsection{Implementation}

The \texttt{VTKModelExporter} was extended with a general cell-data
mechanism:

\begin{enumerate}
  \item A new member \texttt{cell\_fields\_} (same \texttt{FieldBuffer}
        type used for nodal fields).
  \item A new static method \texttt{attach\_cell\_field()} that calls
        \texttt{grid->GetCellData()->AddArray(...)}.
  \item A \texttt{detach\_mesh\_cell\_arrays()} housekeeping method.
  \item Integration into \texttt{write\_mesh()}: after attaching
        point-data arrays, cell-data arrays are also attached.
  \item The \texttt{vtkCellData.h} header was added to the VTK include
        block.
\end{enumerate}

A dedicated public method \texttt{set\_local\_axes(e\_x, e\_y, e\_z)}
creates three 3-component cell fields named \texttt{local\_x},
\texttt{local\_y} and \texttt{local\_z}.  Each field contains one
3D vector per cell, all set to the corresponding axis direction.
Usage:

\begin{lstlisting}[language=C++, basicstyle=\ttfamily\small]
vtk_exporter.set_local_axes(spec.e_x, spec.e_y, spec.e_z);
vtk_exporter.write_mesh("submodel.vtu");
\end{lstlisting}

\subsection{Visualisation in ParaView}

To display local axes in ParaView:

\begin{enumerate}
  \item Open the \texttt{.vtu} file.
  \item Apply a \textbf{Cell Centers} filter to obtain a point per cell.
  \item Apply a \textbf{Glyph} filter on the result:
        \begin{itemize}
          \item Glyph type: Arrow.
          \item Orientation array: \texttt{local\_x} (or \texttt{local\_y},
                \texttt{local\_z}).
          \item Scale mode: Off (uniform arrow length).
        \end{itemize}
  \item Repeat for each axis, using different colours.
\end{enumerate}


% =========================================================================
\section{Example: Hex20 Prismatic Sub-Model}
\label{sec:ho-example}

\begin{lstlisting}[language=C++, basicstyle=\ttfamily\small]
auto spec = fall_n::align_to_beam(
    A, B, up,
    0.40, 0.60,    // width, height
    4, 6, 8,       // nx, ny, nz
    "Concrete",
    fall_n::HexOrder::Serendipity   // Hex20
);

auto [domain, grid] = fall_n::make_prismatic_domain(spec);
// domain contains:
//   4 * 6 * 8 = 192  Hex20 elements  (27 GPs each)
//   (2*4+1)*(2*6+1)*(2*8+1) = 9 * 13 * 17 = 1989 nodes
\end{lstlisting}


% =========================================================================
\section{Backward Compatibility}
\label{sec:ho-compat}

The default value \texttt{HexOrder::Linear} ensures that every existing
call to \texttt{make\_prismatic\_domain()},
\texttt{make\_reinforced\_prismatic\_domain()} and
\texttt{align\_to\_beam()} continues to produce the same Hex8 meshes
without modification.  The grid's \texttt{step = 1} reproduces the
original \texttt{(n+1)} indexing identically.
